\chapter{Assigned Batch Parameters}

\section{Application Scenario}
This project considers an industrial sensor monitoring communication system in which a \\remote sensor node measures process variables 
such as temperature, vibration, pressure, or \\machine health condition and transmits the data wirelessly to a nearby monitoring receiver. 
The transmitted information is represented in binary form and modulated using Binary Phase Shift Keying (BPSK), 
which is selected because of its simple implementation and reliable \\performance in noisy environments.

For our batch, the communication channel is modeled as a mild low-pass channel followed by additive white Gaussian noise (AWGN). 
At the receiver side, both the correlator receiver and the FIR matched filter receiver are implemented and compared in order 
to study their detection performance under the assigned channel condition.

\section{Numerical Specifications}
The batch-specific numerical parameters assigned for this project are listed in Table~\ref{tab:batch9params}. 
These values are used throughout the Python implementation, simulation, and analysis.

\begin{table}[H]
\centering
\caption{Assigned numerical parameters}
\label{tab:batch9params}
\begin{tabular}{|l|l|}
\hline
\textbf{Parameter} & \textbf{Assigned Value} \\
\hline
Application theme & Industrial Sensor Monitoring Communication \\
\hline
Carrier frequency, \(f_c\) & 15 kHz \\
\hline
Bit rate, \(R_b\) & 1 kbps \\
\hline
Bit duration, \(T_b = 1/R_b\) & 1 ms \\
\hline
Sampling frequency, \(f_s\) & 200 kHz \\
\hline
Samples per bit, \(N = f_s/R_b\) & 200 \\
\hline
Signal amplitude, \(A\) & 1 V \\
\hline
Number of bits for BER simulation & \(10^5\) \\
\hline
\(E_b/N_0\) range & 2 dB to 14 dB \\
\hline
Channel type & Mild LPF + AWGN \\
\hline
Low-pass channel impulse response & \(h[n] = [0.2,\;0.6,\;0.2]\) \\
\hline
Receiver type & Both correlator and FIR matched filter \\
\hline
Decision threshold & 0 \\
\hline
\end{tabular}
\end{table}

\section{Derived Quantities}
From the assigned numerical values, a few important derived quantities are calculated and used in later chapters. 
These quantities help in waveform generation, spectrum analysis, and receiver implementation.

The bit duration is
\[
T_b = \frac{1}{R_b} = \frac{1}{1000} = 0.001~\text{s} = 1~\text{ms}
\]

The number of samples per bit is
\[
N = \frac{f_s}{R_b} = \frac{200000}{1000} = 200
\]

The theoretical energy per bit for BPSK is
\[
E_b = \frac{A^2 T_b}{2}
\]
Since \(A = 1\) V and \(T_b = 0.001\) s,
\[
E_b = \frac{(1)^2 \times 0.001}{2} = 0.0005~\text{J}
\]

The null-to-null bandwidth of BPSK is
\[
BW = 2R_b = 2 \times 1000 = 2000~\text{Hz}
\]

The first spectral null frequencies are
\[
f_{\text{lower null}} = f_c - R_b = 15000 - 1000 = 14000~\text{Hz}
\]
\[
f_{\text{upper null}} = f_c + R_b = 15000 + 1000 = 16000~\text{Hz}
\]

These values are important because they define the expected frequency-domain shape of the transmitted BPSK signal and help validate the FFT and PSD plots.

\section{Discussion}
The assigned parameter set corresponds to a practical low-data-rate industrial monitoring link, where reliable transmission is more important than very high speed. 
The presence of a mild low-pass channel makes the problem more realistic because real industrial links often introduce slight distortion in addition to random noise.

Since both correlator and matched filter receivers are assigned, this batch provides a good \\opportunity to compare the two detection methods and verify their equivalence in decision \\making. 
The BER and capacity results obtained later in the report will therefore reflect both \\theoretical understanding and practical simulation performance.